\documentclass[11pt, a4paper]{article}
\usepackage[a4paper, top=2.5cm, bottom=2.5cm, left=2cm, right=2cm]{geometry}
\usepackage{fontspec}
\usepackage[english, bidi=basic, provide=*]{babel}
\babelprovide[import, onchar=ids fonts]{english}
\babelfont{rm}{TeX Gyre Termes}
\babelfont{sf}{TeX Gyre Heros}
\babelfont{tt}{TeX Gyre Cursor}
\usepackage{enumitem}
\usepackage{amsmath}

\title{\textbf{Chapter 2: The Chemical Context of Life \\ 75 Multiple Choice Questions}}
\author{}
\date{}

\begin{document}

\maketitle

\begin{enumerate}

% ---------------------------------------
% 2.1 Elements and Compounds (1-10)
% ---------------------------------------

\item \textbf{Which of the following is defined as anything that takes up space and has mass?}
\begin{enumerate}[label=\Alph*.]
    \item Energy
    \item Element
    \item Compound
    \item Matter
    \item Molecule
\end{enumerate}
\textbf{Answer:} D \\
\textbf{Explanation:} Matter is the physical substance that makes up the universe, defined by having mass and occupying volume.

\item \textbf{A substance that cannot be broken down to other substances by chemical reactions is known as:}
\begin{enumerate}[label=\Alph*.]
    \item A compound
    \item An element
    \item An isotope
    \item A molecule
    \item A proton
\end{enumerate}
\textbf{Answer:} B \\
\textbf{Explanation:} Elements are the purest forms of matter that cannot be chemically simplified. There are 92 naturally occurring elements.

\item \textbf{Which of the following represents a compound?}
\begin{enumerate}[label=\Alph*.]
    \item $O_2$
    \item $H_2$
    \item C
    \item NaCl
    \item Au
\end{enumerate}
\textbf{Answer:} D \\
\textbf{Explanation:} A compound consists of two or more different elements combined in a fixed ratio. NaCl consists of Sodium and Chlorine.

\item \textbf{Sodium is a highly reactive metal and chlorine is a poisonous gas. When combined, they form edible table salt (NaCl). This illustrates the concept of:}
\begin{enumerate}[label=\Alph*.]
    \item Isotopic decay
    \item Emergent properties
    \item Covalent bonding
    \item Van der Waals interactions
    \item Radiometric dating
\end{enumerate}
\textbf{Answer:} B \\
\textbf{Explanation:} Emergent properties appear at higher levels of organization, where compounds exhibit traits entirely different from their individual constituent elements.

\item \textbf{Which four elements make up approximately 96\% of living matter?}
\begin{enumerate}[label=\Alph*.]
    \item Carbon, Hydrogen, Nitrogen, Oxygen
    \item Carbon, Oxygen, Phosphorus, Sulfur
    \item Calcium, Potassium, Sodium, Chlorine
    \item Carbon, Sodium, Chlorine, Oxygen
    \item Nitrogen, Oxygen, Iron, Iodine
\end{enumerate}
\textbf{Answer:} A \\
\textbf{Explanation:} O, C, H, and N are the foundational elements of the biological macromolecules that comprise all life.

\item \textbf{Which of the following elements is considered a trace element in humans?}
\begin{enumerate}[label=\Alph*.]
    \item Carbon
    \item Oxygen
    \item Hydrogen
    \item Iodine
    \item Nitrogen
\end{enumerate}
\textbf{Answer:} D \\
\textbf{Explanation:} Trace elements are required in minute quantities. Iodine is needed for thyroid hormone production, and deficiency leads to goiter.

\item \textbf{An iodine deficiency in the human diet can lead to an abnormally enlarged thyroid gland, a condition known as:}
\begin{enumerate}[label=\Alph*.]
    \item Anemia
    \item Goiter
    \item Diabetes
    \item Malaria
    \item Scurvy
\end{enumerate}
\textbf{Answer:} B \\
\textbf{Explanation:} A daily intake of only 0.15 mg of iodine is adequate to prevent goiter.

\item \textbf{Which naturally occurring toxic element has been linked to numerous diseases due to contaminated groundwater in southern Asia?}
\begin{enumerate}[label=\Alph*.]
    \item Iron
    \item Calcium
    \item Arsenic
    \item Potassium
    \item Phosphorus
\end{enumerate}
\textbf{Answer:} C \\
\textbf{Explanation:} Arsenic is naturally occurring and toxic, causing severe health issues when consumed in contaminated well water.

\item \textbf{Serpentine plant communities exist in soils containing toxic levels of chromium, nickel, and cobalt. The survival of these plants is best explained by:}
\begin{enumerate}[label=\Alph*.]
    \item Emergent properties
    \item Natural selection and evolutionary adaptation
    \item Radiometric decay
    \item Covalent bonding of toxic elements
    \item Chemical equilibrium
\end{enumerate}
\textbf{Answer:} B \\
\textbf{Explanation:} Variants of ancestral species adapted to the toxic elements, and natural selection led to the specialized serpentine plant communities.

\item \textbf{Essential elements are defined as those that:}
\begin{enumerate}[label=\Alph*.]
    \item Are found in the Earth's crust in large amounts.
    \item An organism needs to live a healthy life and reproduce.
    \item Are only required by plants, not animals.
    \item Make up exactly 100\% of living matter.
    \item Are biologically inert.
\end{enumerate}
\textbf{Answer:} B \\
\textbf{Explanation:} Organisms require essential elements for survival and reproduction. Humans need about 25, while plants need 17.

% ---------------------------------------
% 2.2 Atomic Structure (11-48)
% ---------------------------------------

\item \textbf{The smallest unit of matter that still retains the properties of an element is the:}
\begin{enumerate}[label=\Alph*.]
    \item Proton
    \item Molecule
    \item Atom
    \item Electron
    \item Neutron
\end{enumerate}
\textbf{Answer:} C \\
\textbf{Explanation:} Atoms are the fundamental building blocks of elements.

\item \textbf{Which two subatomic particles are packed tightly together in the atomic nucleus?}
\begin{enumerate}[label=\Alph*.]
    \item Electrons and protons
    \item Neutrons and electrons
    \item Protons and neutrons
    \item Positrons and electrons
    \item Photons and protons
\end{enumerate}
\textbf{Answer:} C \\
\textbf{Explanation:} Protons and neutrons form the dense core (nucleus) of the atom, while electrons orbit in a cloud around it.

\item \textbf{What is the electrical charge of a proton?}
\begin{enumerate}[label=\Alph*.]
    \item Neutral (0)
    \item One unit of positive charge (+1)
    \item One unit of negative charge (-1)
    \item Two units of positive charge (+2)
    \item It depends on the element
\end{enumerate}
\textbf{Answer:} B \\
\textbf{Explanation:} Each proton carries one unit of positive charge.

\item \textbf{The unit of measurement used for the mass of subatomic particles and atoms is the:}
\begin{enumerate}[label=\Alph*.]
    \item Gram
    \item Newton
    \item Joule
    \item Dalton (amu)
    \item Mole
\end{enumerate}
\textbf{Answer:} D \\
\textbf{Explanation:} The dalton (or atomic mass unit, amu) is used because subatomic particles are incredibly small (protons/neutrons are $\sim$1 dalton each).

\item \textbf{Which subatomic particle's mass is considered negligible when computing the total mass of an atom?}
\begin{enumerate}[label=\Alph*.]
    \item Proton
    \item Neutron
    \item Nucleus
    \item Electron
    \item Positron
\end{enumerate}
\textbf{Answer:} D \\
\textbf{Explanation:} The mass of an electron is only about 1/2000th that of a proton or neutron, making its contribution to atomic mass negligible.

\item \textbf{The atomic number of an element indicates the number of \_\_\_\_\_\_\_ in its atoms.}
\begin{enumerate}[label=\Alph*.]
    \item Neutrons
    \item Protons
    \item Electrons
    \item Protons plus neutrons
    \item Valence electrons
\end{enumerate}
\textbf{Answer:} B \\
\textbf{Explanation:} The atomic number is unique to each element and specifies the exact number of protons in the nucleus.

\item \textbf{In an electrically neutral atom, the number of protons is always equal to the number of:}
\begin{enumerate}[label=\Alph*.]
    \item Neutrons
    \item Valence electrons only
    \item Electrons
    \item Orbitals
    \item Daltons
\end{enumerate}
\textbf{Answer:} C \\
\textbf{Explanation:} To maintain a neutral electrical charge, the positively charged protons must be perfectly balanced by the negatively charged electrons.

\item \textbf{The mass number of an atom is defined as:}
\begin{enumerate}[label=\Alph*.]
    \item The total number of protons and electrons.
    \item The total number of protons and neutrons in the nucleus.
    \item The total number of neutrons and electrons.
    \item The atomic mass in grams.
    \item The number of valence electrons.
\end{enumerate}
\textbf{Answer:} B \\
\textbf{Explanation:} Since protons and neutrons each have a mass of roughly 1 dalton, their sum constitutes the mass number.

\item \textbf{An atom of Sodium (Na) has an atomic number of 11 and a mass number of 23. How many neutrons does it have?}
\begin{enumerate}[label=\Alph*.]
    \item 11
    \item 23
    \item 12
    \item 34
    \item 10
\end{enumerate}
\textbf{Answer:} C \\
\textbf{Explanation:} Number of neutrons = Mass number $-$ Atomic number. $23 - 11 = 12$.

\item \textbf{The simplest atom, Hydrogen ($^1H$), differs from most other atoms because it lacks:}
\begin{enumerate}[label=\Alph*.]
    \item Protons
    \item Electrons
    \item Neutrons
    \item A nucleus
    \item Valence shells
\end{enumerate}
\textbf{Answer:} C \\
\textbf{Explanation:} The most common isotope of hydrogen consists of a single proton and a single electron, with no neutrons.

\item \textbf{Different atomic forms of the same element that have the same number of protons but different numbers of neutrons are called:}
\begin{enumerate}[label=\Alph*.]
    \item Isomers
    \item Ions
    \item Isotopes
    \item Compounds
    \item Polymers
\end{enumerate}
\textbf{Answer:} C \\
\textbf{Explanation:} Isotopes vary in their neutron count and therefore in mass, but they retain the same chemical identity.

\item \textbf{Carbon-12 ($^{12}C$) and Carbon-14 ($^{14}C$) behave identically in chemical reactions because they have the same number of:}
\begin{enumerate}[label=\Alph*.]
    \item Neutrons
    \item Protons and Electrons
    \item Orbitals and Neutrons
    \item Atomic mass
    \item Radioactive decays
\end{enumerate}
\textbf{Answer:} B \\
\textbf{Explanation:} Chemical behavior is dictated by electrons (specifically valence electrons), which are determined by the number of protons.

\item \textbf{What happens when a radioactive isotope undergoes decay?}
\begin{enumerate}[label=\Alph*.]
    \item The nucleus spontaneously gives off particles and energy.
    \item The atom gains electrons to become an anion.
    \item The atom chemically bonds with another atom.
    \item The atom changes state from a solid to a gas.
    \item The atom permanently freezes.
\end{enumerate}
\textbf{Answer:} A \\
\textbf{Explanation:} Radioactive decay is the spontaneous loss of subatomic particles and energy from an unstable nucleus.

\item \textbf{When Carbon-14 decays, a neutron transforms into a proton. This changes the atom into:}
\begin{enumerate}[label=\Alph*.]
    \item Carbon-13
    \item Oxygen-16
    \item Nitrogen-14
    \item Hydrogen-1
    \item Helium-4
\end{enumerate}
\textbf{Answer:} C \\
\textbf{Explanation:} Gaining a proton changes the atomic number from 6 to 7, transforming the carbon atom into a nitrogen atom.

\item \textbf{In medical diagnostics, radioactive isotopes are frequently used as:}
\begin{enumerate}[label=\Alph*.]
    \item Antibiotics
    \item Painkillers
    \item Tracers to track atoms during metabolism
    \item Sedatives
    \item Gene editing tools
\end{enumerate}
\textbf{Answer:} C \\
\textbf{Explanation:} Cells incorporate radioactive isotopes identically to stable ones, allowing tracking with imaging instruments like PET scanners.

\item \textbf{A PET scan can detect cancerous tissues by looking for elevated levels of radioactively labeled:}
\begin{enumerate}[label=\Alph*.]
    \item Glucose
    \item Oxygen
    \item Sodium
    \item Water
    \item Hemoglobin
\end{enumerate}
\textbf{Answer:} A \\
\textbf{Explanation:} Cancerous tissues have high metabolic activity and consume large amounts of glucose, which lights up on a PET scan.

\item \textbf{The time it takes for 50\% of a parent radioactive isotope to decay is called its:}
\begin{enumerate}[label=\Alph*.]
    \item Decay rate
    \item Half-life
    \item Atomic duration
    \item Isotopic period
    \item Radiant span
\end{enumerate}
\textbf{Answer:} B \\
\textbf{Explanation:} Half-life is a constant property of radioactive isotopes, unaffected by environmental conditions, and is used for radiometric dating.

\item \textbf{To determine the age of moon rocks or Earth itself, scientists used which long-lived radioactive isotope?}
\begin{enumerate}[label=\Alph*.]
    \item Carbon-14
    \item Nitrogen-14
    \item Uranium-238
    \item Phosphorus-32
    \item Iodine-131
\end{enumerate}
\textbf{Answer:} C \\
\textbf{Explanation:} Uranium-238 has a half-life of 4.5 billion years, making it suitable for dating incredibly ancient rocks.

\item \textbf{Why cannot Carbon-14 dating be used to date dinosaur fossils (which are >65.5 million years old)?}
\begin{enumerate}[label=\Alph*.]
    \item Dinosaurs did not consume carbon.
    \item Carbon-14 only works on plant matter.
    \item Carbon-14's half-life is 5,730 years, so it depletes to undetectable levels long before 65 million years.
    \item Dinosaur bones are fossilized into rock, protecting them from decay.
    \item Carbon-14 turns directly into lead in dinosaur fossils.
\end{enumerate}
\textbf{Answer:} C \\
\textbf{Explanation:} After about 75,000 years, the remaining fraction of C-14 is too infinitesimally small to measure accurately.

\item \textbf{Energy is defined as:}
\begin{enumerate}[label=\Alph*.]
    \item The mass of an object multiplied by the speed of light.
    \item The capacity to cause change, such as doing work.
    \item The sum of all moving atoms in an object.
    \item The amount of heat released by a reaction.
    \item A spontaneous nuclear decay.
\end{enumerate}
\textbf{Answer:} B \\
\textbf{Explanation:} Energy is the ability to cause change or do work, shifting matter against opposing forces.

\item \textbf{The energy that matter possesses because of its location or structure is called:}
\begin{enumerate}[label=\Alph*.]
    \item Kinetic energy
    \item Thermal energy
    \item Activation energy
    \item Potential energy
    \item Chemical equilibrium
\end{enumerate}
\textbf{Answer:} D \\
\textbf{Explanation:} Potential energy is stored energy. For electrons, it is based on their distance from the nucleus.

\item \textbf{Which subatomic particles are directly involved in chemical reactions?}
\begin{enumerate}[label=\Alph*.]
    \item Protons
    \item Neutrons
    \item Electrons
    \item Positrons
    \item Quarks
\end{enumerate}
\textbf{Answer:} C \\
\textbf{Explanation:} Atoms interact via their outer electron clouds; nuclei do not come close enough to interact during standard chemical reactions.

\item \textbf{Regarding the potential energy of electrons, which statement is true?}
\begin{enumerate}[label=\Alph*.]
    \item Electrons closer to the nucleus have higher potential energy.
    \item The potential energy of electrons is constant regardless of distance.
    \item The more distant an electron is from the nucleus, the greater its potential energy.
    \item Electrons have zero potential energy since they have negligible mass.
    \item Electrons absorb energy when moving closer to the nucleus.
\end{enumerate}
\textbf{Answer:} C \\
\textbf{Explanation:} Because it takes work to move an electron away from the positive pull of the nucleus, farther electrons have greater potential energy.

\item \textbf{An electron moves from the third shell to the first shell. What happens to energy during this process?}
\begin{enumerate}[label=\Alph*.]
    \item Energy is absorbed.
    \item Energy is released, usually as light or heat.
    \item Potential energy increases.
    \item The atomic number changes.
    \item A proton is gained.
\end{enumerate}
\textbf{Answer:} B \\
\textbf{Explanation:} Moving to a shell closer to the nucleus implies a drop to a lower energy state. The "lost" energy is released to the environment.

\item \textbf{The chemical behavior of an atom is determined mostly by:}
\begin{enumerate}[label=\Alph*.]
    \item The number of neutrons.
    \item The total mass of the atom.
    \item The number of electrons in its innermost shell.
    \item The number of electrons in its outermost (valence) shell.
    \item The number of protons.
\end{enumerate}
\textbf{Answer:} D \\
\textbf{Explanation:} Valence electrons are the outermost electrons and dictate how an atom interacts and bonds with other atoms.

\item \textbf{What is the maximum number of electrons that the first electron shell can hold?}
\begin{enumerate}[label=\Alph*.]
    \item 2
    \item 4
    \item 8
    \item 10
    \item 18
\end{enumerate}
\textbf{Answer:} A \\
\textbf{Explanation:} The first electron shell contains a single $1s$ orbital, accommodating a maximum of 2 electrons.

\item \textbf{An atom of Fluorine has 7 valence electrons, and an atom of Chlorine has 7 valence electrons. Consequently, they:}
\begin{enumerate}[label=\Alph*.]
    \item Belong to different periods but have opposite chemical behaviors.
    \item Exhibit similar chemical behaviors.
    \item Will both form cations easily.
    \item Are both completely unreactive.
    \item Form triple covalent bonds.
\end{enumerate}
\textbf{Answer:} B \\
\textbf{Explanation:} Elements with the same number of valence electrons exhibit similar chemical properties and bond similarly.

\item \textbf{Elements like Helium, Neon, and Argon possess full valence shells. Therefore, they are:}
\begin{enumerate}[label=\Alph*.]
    \item Highly reactive metals
    \item Isotopes of each other
    \item Inert (chemically unreactive)
    \item Highly radioactive
    \item Strong acids
\end{enumerate}
\textbf{Answer:} C \\
\textbf{Explanation:} A full valence shell means the atom is stable and does not readily interact or form bonds with other atoms.

\item \textbf{The three-dimensional space where an electron is found 90\% of the time is called an:}
\begin{enumerate}[label=\Alph*.]
    \item Electron shell
    \item Atomic nucleus
    \item Orbital
    \item Isotope ring
    \item Electron lattice
\end{enumerate}
\textbf{Answer:} C \\
\textbf{Explanation:} Orbitals represent the probability volumes describing where electrons spend most of their time.

\item \textbf{What is the shape of a $p$ orbital?}
\begin{enumerate}[label=\Alph*.]
    \item Spherical
    \item Dumbbell-shaped
    \item Tetrahedral
    \item Cuboid
    \item Linear
\end{enumerate}
\textbf{Answer:} B \\
\textbf{Explanation:} The $p$ orbitals have a characteristic dumbbell shape, with three $p$ orbitals oriented along the x, y, and z axes.

\item \textbf{The second electron shell holds a maximum of 8 electrons. How many total orbitals does the second shell contain?}
\begin{enumerate}[label=\Alph*.]
    \item One $s$ orbital and one $p$ orbital (2 total)
    \item Four $s$ orbitals (4 total)
    \item One $s$ orbital and three $p$ orbitals (4 total)
    \item Eight orbitals (8 total)
    \item Three $s$ orbitals and one $p$ orbital (4 total)
\end{enumerate}
\textbf{Answer:} C \\
\textbf{Explanation:} The second shell consists of one $2s$ orbital and three $2p$ orbitals. Each of the four orbitals holds 2 electrons ($4 \times 2 = 8$).

% ---------------------------------------
% 2.3 Chemical Bonding (49-68)
% ---------------------------------------

\item \textbf{A covalent bond is formed by:}
\begin{enumerate}[label=\Alph*.]
    \item The complete transfer of electrons from one atom to another.
    \item The sharing of a pair of valence electrons by two atoms.
    \item The attraction between transient partial charges.
    \item The magnetic pull of the nucleus on another atom's electrons.
    \item The loss of neutrons.
\end{enumerate}
\textbf{Answer:} B \\
\textbf{Explanation:} Covalent bonds involve atoms sharing valence electrons to complete their respective valence shells.

\item \textbf{Two or more atoms held together by covalent bonds constitute a:}
\begin{enumerate}[label=\Alph*.]
    \item Salt
    \item Mixture
    \item Solution
    \item Molecule
    \item Cation
\end{enumerate}
\textbf{Answer:} D \\
\textbf{Explanation:} By definition, a molecule consists of multiple atoms bonded together covalently.

\item \textbf{In a double covalent bond, such as in $O_2$, how many pairs of electrons are shared?}
\begin{enumerate}[label=\Alph*.]
    \item One
    \item Two
    \item Three
    \item Four
    \item None
\end{enumerate}
\textbf{Answer:} B \\
\textbf{Explanation:} A double bond involves sharing two pairs of valence electrons (a total of 4 electrons).

\item \textbf{Carbon has 4 valence electrons. What is its bonding capacity (valence)?}
\begin{enumerate}[label=\Alph*.]
    \item 1
    \item 2
    \item 3
    \item 4
    \item 8
\end{enumerate}
\textbf{Answer:} D \\
\textbf{Explanation:} Carbon needs 4 additional electrons to fill its valence shell (which holds 8), so it can form 4 covalent bonds.

\item \textbf{The attraction of a particular atom for the electrons of a covalent bond is called:}
\begin{enumerate}[label=\Alph*.]
    \item Polarity
    \item Radioactivity
    \item Electronegativity
    \item Ionic potential
    \item Hybridization
\end{enumerate}
\textbf{Answer:} C \\
\textbf{Explanation:} Electronegativity measures how strongly an atom pulls shared electrons toward its own nucleus.

\item \textbf{Why is the covalent bond in a hydrogen molecule ($H_2$) considered nonpolar?}
\begin{enumerate}[label=\Alph*.]
    \item Hydrogen has no electronegativity.
    \item The electrons are stripped completely by one atom.
    \item Both atoms have identical electronegativity, so electrons are shared equally.
    \item Hydrogen atoms only have one electron shell.
    \item It relies on van der Waals interactions.
\end{enumerate}
\textbf{Answer:} C \\
\textbf{Explanation:} When atoms of the same element bond, their electronegativities are equal, resulting in a standoff and equal electron sharing.

\item \textbf{In a water molecule ($H_2O$), the oxygen atom is more electronegative than the hydrogen atoms. As a result:}
\begin{enumerate}[label=\Alph*.]
    \item Oxygen strips electrons completely to form $O^{2-}$.
    \item The bond is nonpolar covalent.
    \item Electrons spend more time near the oxygen, creating a partial negative charge on oxygen.
    \item Electrons spend more time near hydrogen, creating a partial positive charge on oxygen.
    \item The molecule forms an ionic lattice.
\end{enumerate}
\textbf{Answer:} C \\
\textbf{Explanation:} Unequal sharing of electrons toward the more electronegative oxygen makes the bonds polar covalent, yielding partial charges ($\delta-$ on O, $\delta+$ on H).

\item \textbf{When two atoms are so unequal in electronegativity that one strips an electron completely away from the other, they form:}
\begin{enumerate}[label=\Alph*.]
    \item Polar covalent bonds
    \item Nonpolar covalent bonds
    \item Ions and ionic bonds
    \item Hydrogen bonds
    \item Isotopes
\end{enumerate}
\textbf{Answer:} C \\
\textbf{Explanation:} This complete transfer results in oppositely charged ions, which then attract each other to form an ionic bond.

\item \textbf{A positively charged ion is called a(n) \_\_\_\_\_\_\_\_, and a negatively charged ion is called a(n) \_\_\_\_\_\_\_\_.}
\begin{enumerate}[label=\Alph*.]
    \item Anion; Cation
    \item Cation; Anion
    \item Proton; Electron
    \item Positron; Negatron
    \item Isotope; Isomer
\end{enumerate}
\textbf{Answer:} B \\
\textbf{Explanation:} Cations are positive (think of the 't' as a plus sign), and anions are negative (a-negative-ion).

\item \textbf{Sodium chloride (NaCl) is an example of an ionic compound, commonly called a salt. Which statement about NaCl is true?}
\begin{enumerate}[label=\Alph*.]
    \item It consists of distinct, individual NaCl molecules.
    \item It consists of a 3D lattice of cations and anions.
    \item It shares electrons covalently.
    \item It is virtually unbreakable even when dissolved in water.
    \item It acts as a single large isotope.
\end{enumerate}
\textbf{Answer:} B \\
\textbf{Explanation:} Salts do not consist of molecules; they are vast crystalline aggregates of ions held by electrical attraction.

\item \textbf{How does water affect ionic bonds?}
\begin{enumerate}[label=\Alph*.]
    \item Water strengthens ionic bonds.
    \item Water converts ionic bonds into covalent bonds.
    \item Water weakens ionic bonds because each ion becomes shielded by water molecules.
    \item Water evaporates, leaving pure sodium metal.
    \item Water reacts with salts to form toxic gases.
\end{enumerate}
\textbf{Answer:} C \\
\textbf{Explanation:} In aqueous solutions, water molecules surround and separate the ions, making the ionic bonds much weaker than in dry crystals.

\item \textbf{A hydrogen bond is defined as a noncovalent attraction between:}
\begin{enumerate}[label=\Alph*.]
    \item Two hydrogen atoms.
    \item A hydrogen atom (with a partial positive charge) and another electronegative atom (with a partial negative charge).
    \item Two atoms with identical electronegativities.
    \item Any cation and anion.
    \item Two completely nonpolar molecules.
\end{enumerate}
\textbf{Answer:} B \\
\textbf{Explanation:} Hydrogen bonds form when a $\delta+$ hydrogen is attracted to a nearby $\delta-$ atom (like Oxygen or Nitrogen) on another molecule.

\item \textbf{Gecko lizards can walk straight up a wall due to which type of weak interactions?}
\begin{enumerate}[label=\Alph*.]
    \item Covalent bonds
    \item Ionic bonds
    \item Van der Waals interactions
    \item Hydrogen bonds
    \item Radioactive forces
\end{enumerate}
\textbf{Answer:} C \\
\textbf{Explanation:} The ever-changing regions of positive and negative charge (van der Waals interactions) across millions of tiny hairs provide enough cumulative attraction to support the gecko.

\item \textbf{When an atom forms covalent bonds, its $s$ and $p$ orbitals may hybridize. In carbon (e.g., $CH_4$), these four hybrid orbitals point toward the corners of a:}
\begin{enumerate}[label=\Alph*.]
    \item Square
    \item Triangle
    \item Tetrahedron
    \item Hexagon
    \item Sphere
\end{enumerate}
\textbf{Answer:} C \\
\textbf{Explanation:} The hybridization of one $s$ and three $p$ orbitals creates four teardrop-shaped orbitals forming a tetrahedron (pyramid with a triangular base).

\item \textbf{Endorphins are naturally produced by the brain to relieve pain. Opiates like morphine have similar effects because:}
\begin{enumerate}[label=\Alph*.]
    \item They covalently destroy pain receptors.
    \item They are identical in chemical formula to endorphins.
    \item They have a molecular shape similar to endorphins and can bind to the same brain receptors.
    \item They undergo rapid radioactive decay in the brain.
    \item They alter the genetic code of the brain cells.
\end{enumerate}
\textbf{Answer:} C \\
\textbf{Explanation:} Molecular shape is crucial; the complementary shapes allow morphine to mimic natural endorphins and bind to specific cellular receptors.

% ---------------------------------------
% 2.4 Chemical Reactions (69-75)
% ---------------------------------------

\item \textbf{The making and breaking of chemical bonds to change the composition of matter is known as:}
\begin{enumerate}[label=\Alph*.]
    \item Chemical equilibrium
    \item Isotopic decay
    \item A chemical reaction
    \item Hybridization
    \item Van der Waals interaction
\end{enumerate}
\textbf{Answer:} C \\
\textbf{Explanation:} Chemical reactions involve rearranging matter by making and breaking bonds between atoms.

\item \textbf{In the chemical reaction $2H_2 + O_2 \rightarrow 2H_2O$, the starting materials ($H_2$ and $O_2$) are called the:}
\begin{enumerate}[label=\Alph*.]
    \item Products
    \item Isotopes
    \item Reactants
    \item Hybrids
    \item Ions
\end{enumerate}
\textbf{Answer:} C \\
\textbf{Explanation:} Reactants are the starting materials on the left side of the reaction arrow.

\item \textbf{According to the principle of conservation of matter in chemical reactions:}
\begin{enumerate}[label=\Alph*.]
    \item Atoms can be created from pure energy.
    \item Atoms are destroyed to release heat.
    \item Matter is rearranged; atoms are neither created nor destroyed.
    \item Total molecular volume is always strictly conserved.
    \item All reactants are converted permanently to products without possibility of reversal.
\end{enumerate}
\textbf{Answer:} C \\
\textbf{Explanation:} Chemical reactions simply rearrange electrons and bonds. The exact number and types of atoms are conserved.

\item \textbf{In the overall equation for photosynthesis ($6CO_2 + 6H_2O \rightarrow C_6H_{12}O_6 + 6O_2$), what provides the energy required to rearrange the matter?}
\begin{enumerate}[label=\Alph*.]
    \item Heat from the soil
    \item Radioactive decay of Carbon-14
    \item Sunlight
    \item Kinetic energy of water
    \item Oxygen gas
\end{enumerate}
\textbf{Answer:} C \\
\textbf{Explanation:} Sunlight powers the conversion of raw ingredients ($CO_2$ and $H_2O$) into glucose and oxygen.

\item \textbf{The point at which the forward and reverse reactions occur at the exact same rate is called:}
\begin{enumerate}[label=\Alph*.]
    \item Reaction termination
    \item Chemical stoichiometry
    \item Chemical equilibrium
    \item The point of no return
    \item Kinetic freeze
\end{enumerate}
\textbf{Answer:} C \\
\textbf{Explanation:} Chemical equilibrium is dynamic; reactions continue in both directions at equal rates, so there is no net change in concentrations.

\item \textbf{At chemical equilibrium, the concentrations of reactants and products are:}
\begin{enumerate}[label=\Alph*.]
    \item Always perfectly equal to each other.
    \item Decreasing constantly.
    \item Stabilized at a particular ratio, and stop changing.
    \item Zero, as the reaction has stopped.
    \item Increasing exponentially.
\end{enumerate}
\textbf{Answer:} C \\
\textbf{Explanation:} Equilibrium does not mean reactants and products are equal in concentration, just that their concentrations have stopped changing because the forward and reverse rates offset.

\item \textbf{If a chemical reaction is indicated by two opposite-headed arrows ($\rightleftharpoons$), it implies that the reaction is:}
\begin{enumerate}[label=\Alph*.]
    \item Highly explosive
    \item Reversible
    \item Occurring in a vacuum
    \item Radioactively decaying
    \item An ionic dissociation
\end{enumerate}
\textbf{Answer:} B \\
\textbf{Explanation:} The double arrow notation indicates that products can re-form reactants, making the reaction reversible.

\end{enumerate}

\end{document}
